\documentclass[11pt,a4paper]{article}
\usepackage[utf8]{inputenc}
\usepackage[spanish]{babel}
\usepackage{graphicx} % Necesario si vas a descomentar la imagen

\begin{document}

\section{Arquitectura}

Para la implementación del sistema, se ha optado por el uso de una Arquitectura desacoplada de tres niveles. Esta se divide en: nivel de presentación siendo este la interfaz de usuario que estará plasmada en una aplicación Android, nivel de aplicación con la lógica de negocio para el Backend API, y, nivel de datos que contendrá la información de los medicamentos almacenados en una base de datos relacional\cite{1}.

El motivo por el cual se está utilizando este tipo de arquitectura es por la flexibilidad que brinda a la hora de separar la lógica del frontend y backend del sistema, permitiendo la implementación de cada una en diferentes lenguajes de programación siendo así más sencillo porque permite que el desarrollo de cada nivel sea simultáneo y asegura la escalabilidad de cada uno sin afectar al resto del sistema\cite{1}.

% Si utilizas alguna imagen aquí iría la figura
% \begin{figure}[h]
%     \centering
%     \includegraphics[width=0.8\textwidth]{tu_imagen.png}
%     \caption{Arquitectura del Sistema para la Aplicación}
%     \label{fig:arquitectura}
% \end{figure}

\subsection{Nivel de Presentación}

Dentro del nivel de presentación, se disponen los elementos necesarios para la interfaz de usuario de la aplicación:

\begin{itemize}
    \item \textbf{Vistas}: Las pantallas que verán los usuarios durante la navegación dentro de la app, donde tendremos elementos como la pantalla principal de la aplicación con el buscador, la información y equivalencias de un medicamento, etc.
    \item \textbf{ViewModel}: Capa intermedia entre interfaz y lógica de app haciendo de intermediario entre su comunicación. 
    \item \textbf{Repository}: Actúa como la única fuente de verdad (\textit{Single Source of Truth}) para la obtención de datos en la aplicación. Su objetivo principal es abstraer el origen de los datos a las capas superiores (\textit{ViewModel}), resolviendo y orquestando las peticiones de red hacia el servidor y abstrayendo la complejidad técnica del consumo de la API o el almacenamiento local.
\end{itemize}

\subsection{Nivel de Aplicación}

Este nivel alberga la lógica de negocio central de la solución y opera como intermediario entre la aplicación móvil y el modelo de datos. Está compuesto por el Backend API, el cual expone y procesa las funcionalidades del sistema mediante los siguientes componentes:

\begin{itemize}
    \item \textbf{API REST}: Funciona como el punto de entrada de todas las comunicaciones externas. Expone diversos \textit{endpoints} HTTP (rutas web) que permiten al nivel de presentación enviar peticiones (por ejemplo, búsquedas por principio activo o peticiones de equivalencias) y recibir una respuesta estructurada en formato JSON.
    \item \textbf{Servicios de Datos}: Capa intermedia dentro del Backend que asume la responsabilidad de procesar y aplicar las reglas de negocio pertinentes. Estos servicios interpretan lo que llega desde la API REST, interactúan con la capa de persistencia para extraer la información y preparan/formatean los registros en objetos de respuesta de forma limpia y manejable para el cliente.
    \item \textbf{Servicio ETL (Extract, Transform, Load)}: Módulo integral encargado de la ingesta de información. Su cometido es conectarse o leer los distintos ficheros crudos de varios países, aplicar transformaciones complejas (limpieza de cadenas, recolección de dosajes y formas farmacéuticas, análisis del código ATC) y cargar finalmente la información de forma unificada y estandarizada en el esquema de la base de datos de producción.
    \item \textbf{Obtención de Equivalencias}: Es el módulo algorítmico responsable de cruzar los medicamentos de un país de origen con los equivalentes de un país destino. Lleva a cabo el emparejamiento basándose primordialmente en una coincidencia exacta de los códigos ATC y, alternativamente, como medida de respaldo, buscando coincidencias en los principios activos del propio fármaco.
\end{itemize}

\subsection{Nivel de Datos}

El nivel de datos se encarga de todo lo relativo al almacenamiento robusto e íntegro de la información, permitiendo la persistencia a largo plazo y la recopilación de registros:

\begin{itemize}
    \item \textbf{Capa de Persistencia}: Materializada en un Sistema de Gestión de Bases de Datos Relacional (como PostgreSQL). Almacena de manera estructurada las distintas entidades como laboratorios, principios activos, medicamentos concretos y, sobre todo, aquellas complejas tablas intermedias que almacenan la relación (N:M). Garantiza las transacciones, un acceso de lectura eficiente para las búsquedas y la integridad referencial.
    \item \textbf{Fuentes Externas}: Conforman los registros en crudo y repositorios ajenos al sistema de los cuales se parte para poblar la capa de persistencia mediante el ETL. Principalmente, están compuestas por listados, \textit{datasets} y bases de datos procedentes de las agencias de salud y reguladoras oficiales correspondientes a cada región soportada por el sistema (tales como la AEMPS española, la FDA americana, el ISP de Chile, etc.).
\end{itemize}
\section{Diseño de Software}

A continuación, se detalla el diseño de las diferentes piezas de software que componen la solución. Debido a la naturaleza modular de esta aplicación, este diseño se ha dividido en distintos bloques lógicos que abarcan la ingesta de información, la gestión de persistencia y las funcionalidades expuestas hacia la capa de presentación. 

\subsection{Base de Datos}

El diseño de la base de datos es un pilar fundamental en el desarrollo de la aplicación, ya que garantiza la integridad, estructuración y el acceso eficiente a la información médica internacional. Atendiendo al diseño conceptual, se ha elaborado un Modelo Entidad-Relación (MER) que plasma las entidades involucradas y cómo se interrelacionan, tal como se aprecia en la Figura \ref{fig:diagrama_mer}.

\begin{figure}[h]
    \centering
    \includegraphics[width=\textwidth]{assets/Diagrama_MER.png}
    \caption{Diagrama Entidad-Relación (MER) del dominio de datos}
    \label{fig:diagrama_mer}
\end{figure}

El esquema de persistencia (cuya estructura de definición de datos se encuentra plasmada en los \textit{scripts} SQL del directorio \texttt{modulo\_datos/database/}) se divide lógicamente en dos grandes áreas:
\begin{itemize}
    \item \textbf{Esquema de \textit{fuentes}}: Actúa como área de \textit{staging} o zona de aterrizaje transitoria para los datos crudos extraídos directamente de las agencias gubernamentales. Alberga tablas independientes sin relaciones de integridad entre sí para cada país (\texttt{spain\_med}, \texttt{chile\_med}, \texttt{canada\_med}, \texttt{usa\_med}, \texttt{portugal\_med}).
    \item \textbf{Esquema \textit{med}}: Representa el núcleo del modelo relacional normalizado. Aquí es donde convergen de forma estandarizada todos los datos ya limpios y listos para ser consumidos.
\end{itemize}

El paso a tablas relacionales de este modelo, ilustrado en la Figura \ref{fig:tablas_mer}, muestra cómo se implementa físicamente el esquema de producción (\textit{med}), el cual está constituido por tres elementos clave:
\begin{itemize}
    \item \textbf{Tablas Maestras o Diccionarios}: Entidades fuertes y descriptivas que albergan valores estandarizados (\texttt{pais}, \texttt{atc}, \texttt{principio\_activo}, \texttt{forma\_farmaceutica} y \texttt{via\_administracion}).
    \item \textbf{Entidad Principal (\texttt{medicamento})}: Almacena los atributos fundamentales e identificativos inherentes a cada fármaco (nombre comercial, laboratorio, país de registro) y se nutre mediante claves foráneas (FK) de origen, forma y vía correspondiente.
    \item \textbf{Tablas Intermedias de Relación (N:M)}: Estructuras relacionales vitales para soportar la multiplicidad del dominio médico. Destacan la tabla \texttt{contiene} (que vincula un medicamento con uno o varios de sus principios activos permitiendo detallar la dosis numérica y de unidad exacta), \texttt{identificado\_por} (asocia un medicamento con sus códigos técnicos ATC) y \texttt{asociado\_con} (enlaza principios activos genéricos con el estándar resolutivo ATC).
\end{itemize}

\begin{figure}[h]
    \centering
    \includegraphics[width=0.95\textwidth]{assets/Tablas_MER.png}
    \caption{Paso a tablas del Modelo Físico Relacional}
    \label{fig:tablas_mer}
\end{figure}

\subsection{Módulo de ETL}

El módulo de ETL (Extracción, Transformación y Carga) es el componente responsable de centralizar, estandarizar y poblar la base de datos de producción con el catálogo internacional de medicamentos. Su funcionamiento se puede apreciar globalmente a través del diagrama de actividad expuesto en la Figura \ref{fig:diagrama_etl}, el cual divide estructuralmente los tres procesos en carriles (\textit{swimlanes}):

\begin{figure}[h]
    \centering
    \includegraphics[width=\textwidth]{assets/Diagrama_Actividad_ETL.png}
    \caption{Diagrama de Actividad del proceso ETL}
    \label{fig:diagrama_etl}
\end{figure}

\begin{enumerate}
    \item \textbf{Extracción de Datos}: Consiste en la obtención inicial de los datos a través de diversas fuentes externas. Por un lado, se consultan de manera automatizada las APIs oficiales (como CIMA o RxNorm), manejando posibles errores de conexión y reintentos de consulta. En paralelo, y para fuentes que no disponen de este tipo de interfaces (como ISPCh en Chile o Infarmed en Portugal), se obtienen datos procedentes de archivos \textit{CSV}. Toda esta información recién adquirida se extrae y se almacena en estado crudo estableciendo conexión y enviándola a un esquema propio de \textit{fuentes} en PostgreSQL.
    
    \item \textbf{Transformación de Datos}: En este paso se toma la carga de datos en crudo y se evalúa si dichos datos ya se encuentran en el sistema. En caso afirmativo, se omiten duplicados y se actualizan los desfasados. De tratarse de un registro nuevo, se procede a extraer sus entidades únicas (Países, ATCs, Principios Activos, Vías de Administración y Formas Farmacéuticas). Además, si se determina que conforma más de un principio activo, el sistema divide y separa mediante delimitadores para aplicar una limpieza individual de las cadenas de texto; si sólo tiene un principio activo se somete a una limpieza y estandarización más básica. Como cierre, todos estos elementos transformados se mapean estructurando diccionarios maestros en memoria.
    
    \item \textbf{Carga en Base de Datos}: Representa la última etapa de consolidación hacia las tablas relacionales finales. Este proceso de carga se divide físicamente en tres rutinas \textit{scripts} secuenciales implementados en Python dentro del directorio \texttt{modulo\_datos/scripts/etl/}:
    \begin{itemize}
        \item \textbf{\texttt{01\_load\_diccionarios.py}}: Actúa como el primer eslabón efectuando la limpieza de registros previos mediante un proceso de borrado/truncado. Posteriormente, toma las entidades previamente estructuradas y ejecuta la carga sobre las tablas de diccionarios maestros, albergando la información base independiente del medicamento (Laboratorios, Entidades Geográficas, Formas Farmacéuticas, Principios Activos, etc.).
        \item \textbf{\texttt{02\_load\_medicamentos.py}}: Se encarga de volcar el volumen más pesado, realizando la carga e inserción completa de la información básica y principal del propio medicamento desglosado por país (identificador, nombre, ATC, entre otros).
        \item \textbf{\texttt{03\_load\_relaciones.py}}: Por último, se responsabiliza de enlazar las partes, consolidando las relaciones de tipo N:M en base de datos; esto es lo que le permite al sistema cruzar qué medicamentos contienen qué principios activos en determinadas cantidades, posibilitando a su vez, la obtención de equivalencias internacionales.
    \end{itemize}
\end{enumerate}

\end{document}